\documentclass{article}
\usepackage{amsmath,amssymb,geometry}
\geometry{margin=1in}

\begin{document}

\section*{CSE 400 – Fundamentals of Probability in Computing}

\subsection*{Lecture 16 Scribe: Joint Random Variables and Joint Distributions}

Based strictly on Lecture Slides by Dhaval Patel, PhD

\bigskip

\section*{1. Why Study Joint Random Variables?}

In many real-world scenarios, we are interested in analyzing two or more random variables simultaneously instead of studying them independently. Joint analysis allows us to understand relationships, dependencies, and correlations between variables.

Examples of pairs of variables often studied jointly:

\begin{itemize}
\item Height and weight of giraffes
\item IQ and birthweight of children
\item Frequency of exercise and rate of heart disease
\item Air pollution level and respiratory illness rate
\item Number of Facebook friends and age of members
\end{itemize}

These examples motivate the need for \textbf{joint random variables}.

\bigskip

\section*{2. Applications Motivating Joint Random Variables}

\subsection*{Example 1: Smart Transportation}

Two random variables are defined:

\[
X = \text{Vehicle Speed}
\]

\[
Y = \text{Traffic Density}
\]

Question posed:

Can speed be high when the traffic density is also high?

Conclusion: The variables may be correlated.

Applications in Intelligent Transportation Systems (ITS):

\begin{itemize}
\item Adaptive traffic signals
\item Congestion prediction
\item Autonomous vehicle routing
\end{itemize}

\bigskip

\subsection*{Example 2: Wireless Network Performance}

Two important random variables:

\[
X = \text{Signal-to-Noise Ratio (SNR)}
\]

\[
Y = \text{Data Throughput}
\]

Observation:

Throughput depends on SNR. Both values vary due to:

\begin{itemize}
\item Fading
\item Interference
\item Network congestion
\end{itemize}

Question:

If SNR is high, will throughput always be high?

Answer discussed: Not necessarily, because network congestion can affect throughput.

Example probability table:

\bigskip

\begin{tabular}{|c|c|c|}
\hline
SNR Level & High Throughput & Low Throughput \\
\hline
High & 0.40 & 0.10 \\
\hline
Medium & 0.20 & 0.15 \\
\hline
Low & 0.05 & 0.10 \\
\hline
\end{tabular}

\bigskip

Applications:

\begin{itemize}
\item 5G scheduling
\item Adaptive modulation
\item Network planning
\end{itemize}

\bigskip

\subsection*{Additional Examples}

Example 3:

\[
X = \text{Rainfall}
\]

\[
Y = \text{Temperature}
\]

Application: Weather forecasting

Example 4:

\[
X = \text{Wind Speed}
\]

\[
Y = \text{Delivery Time}
\]

Application: Drone delivery and logistics systems

Example 5:

Rolling two dice:

\[
X = \text{Outcome of the first die}
\]

\[
Y = \text{Outcome of the second die}
\]

Example questions:

\begin{itemize}
\item $X + Y = 7$
\item $X > Y$
\end{itemize}

These require the use of \textbf{joint probability distributions}.

\bigskip

\section*{3. Discrete Case – Joint PMF}

For discrete random variables $X$ and $Y$, probabilities are assigned to ordered pairs $(x, y)$.

The \textbf{joint probability mass function (PMF)} is defined as:

\[
p(x, y) = P(X = x, Y = y)
\]

The joint PMF describes the probability that $X$ takes value $x$ and $Y$ takes value $y$ simultaneously.

\bigskip

\section*{4. Joint PMF Table Representation}

A common way to represent a joint PMF is using a \textbf{table}.

Structure:

\begin{itemize}
\item Rows correspond to values of $X$
\item Columns correspond to values of $Y$
\item Each cell contains the probability $p(x, y)$
\end{itemize}

Properties of a joint PMF:

\begin{enumerate}
\item $p(x, y) \geq 0$ for all $x$ and $y$
\item The sum of all probabilities equals 1
\end{enumerate}

Sum over all $x$ and $y$:

\[
\sum \sum p(x, y) = 1
\]

\bigskip

\section*{5. Example: Rolling Two Dice}

Define random variables:

\[
X = \text{value of the first die}
\]

\[
Y = \text{value of the second die}
\]

Possible values:

\[
X \in \{1,2,3,4,5,6\}
\]

\[
Y \in \{1,2,3,4,5,6\}
\]

Total possible ordered outcomes:

\[
36
\]

Each pair $(x, y)$ occurs with probability:

\[
P(X = x, Y = y) = \frac{1}{36}
\]

Example events:

Event 1: $X + Y = 7$

Possible pairs:

\[
(1,6)
\]

\[
(2,5)
\]

\[
(3,4)
\]

\[
(4,3)
\]

\[
(5,2)
\]

\[
(6,1)
\]

Number of favorable outcomes = 6

\[
P(X + Y = 7) = 6 / 36 = 1/6
\]

Event 2: $X > Y$

List of ordered pairs where $X$ is greater than $Y$:

\[
(2,1)
\]

\[
(3,1)\ (3,2)
\]

\[
(4,1)\ (4,2)\ (4,3)
\]

\[
(5,1)\ (5,2)\ (5,3)\ (5,4)
\]

\[
(6,1)\ (6,2)\ (6,3)\ (6,4)\ (6,5)
\]

Total favorable outcomes = 15

\[
P(X > Y) = 15 / 36 = 5 / 12
\]

\bigskip

\section*{6. Continuous Case – Joint PDF}

For continuous random variables $X$ and $Y$, the distribution is described by a \textbf{joint probability density function (PDF)}.

Notation:

\[
f(x, y)
\]

Properties:

\begin{enumerate}
\item $f(x, y) \geq 0$
\item Total probability must equal 1:
\end{enumerate}

\[
\int\int f(x, y) \, dx \, dy = 1
\]

Probability over a region $R$ is computed as:

\[
P((X, Y) \in R) = \int\int_R f(x, y) \, dx \, dy
\]

This can be interpreted geometrically as the \textbf{volume under the surface $f(x,y)$} over the region $R$.

\bigskip

\section*{7. Worked Example (Continuous Joint PDF)}

Given joint PDF:

\[
f(x, y) = 4xy
\]

for

\[
0 < x < 1
\]

\[
0 < y < 1
\]

This defines a joint distribution over the \textbf{unit square}.

Tasks:

\begin{enumerate}
\item Verify that $f(x,y)$ is a valid joint PDF.
\item Visualize the event $A = \{X < 0.5, Y > 0.5\}$.
\item Compute $P(A)$.
\end{enumerate}

\bigskip

\section*{8. Example Solution}

\subsection*{Step 1: Verify normalization}

Compute the double integral over the support region.

\[
\int_{0}^{1} \int_{0}^{1} 4xy \, dx \, dy
\]

First integrate with respect to $x$:

\[
\int_{0}^{1} 4xy \, dx
\]

\[
= 4y \int_{0}^{1} x \, dx
\]

\[
= 4y \left[\frac{x^{2}}{2}\right]_{0}^{1}
\]

\[
= 4y \left(\frac{1}{2}\right)
\]

\[
= 2y
\]

Now integrate with respect to $y$:

\[
\int_{0}^{1} 2y \, dy
\]

\[
= \left[y^{2}\right]_{0}^{1}
\]

\[
= 1
\]

Thus the total probability equals 1, so the function is a valid joint PDF.

\bigskip

\subsection*{Step 2: Compute probability of event A}

Event A:

\[
X < 0.5
\]

\[
Y > 0.5
\]

Probability:

\[
P(A) = \int_{0}^{0.5} \int_{0.5}^{1} 4xy \, dy \, dx
\]

First integrate with respect to $y$:

\[
\int_{0.5}^{1} 4xy \, dy
\]

\[
= 4x \int_{0.5}^{1} y \, dy
\]

\[
= 4x \left[\frac{y^{2}}{2}\right]_{0.5}^{1}
\]

\[
= 4x \left(\frac{1 - 0.25}{2}\right)
\]

\[
= 4x \left(\frac{0.75}{2}\right)
\]

\[
= 4x (0.375)
\]

\[
= 1.5x
\]

Now integrate with respect to $x$:

\[
\int_{0}^{0.5} 1.5x \, dx
\]

\[
= 1.5 \left[\frac{x^{2}}{2}\right]_{0}^{0.5}
\]

\[
= 1.5 \left(\frac{0.25}{2}\right)
\]

\[
= 1.5 (0.125)
\]

\[
= 0.1875
\]

\[
= 3 / 16
\]

Therefore:

\[
P(A) = 3/16
\]

\bigskip

\section*{9. Joint Cumulative Distribution Function (Joint CDF)}

Let $X$ and $Y$ be jointly distributed random variables.

The \textbf{joint cumulative distribution function} is defined as:

\[
F(x, y) = P(X < x, Y < y)
\]

This represents the probability that $X$ is less than $x$ and $Y$ is less than $y$ simultaneously.

\bigskip

\section*{10. Joint CDF – Continuous Case}

If $X$ and $Y$ are continuous with joint density $f(x,y)$, the joint CDF is obtained by integrating the joint PDF.

\[
F(x, y) = \int\int f(u, v) \, du \, dv
\]

The joint PDF can be recovered from the joint CDF using partial derivatives:

\[
f(x, y) = \frac{\partial^{2}F(x, y)}{\partial x \, \partial y}
\]

\bigskip

\section*{11. Joint CDF – Discrete Case}

If $X$ and $Y$ are discrete with joint PMF $p(x, y)$, the joint CDF is computed by summing probabilities.

\[
F(x, y) = P(X \leq x, Y \leq y)
\]

This is obtained by summing:

\[
p(x_i, y_j)
\]

for all pairs satisfying:

\[
x_i \leq x
\]

\[
y_j \leq y
\]

\bigskip

\section*{12. Properties of the Joint CDF}

\begin{enumerate}
\item The joint CDF is defined as:
\end{enumerate}

\[
F(x, y) = P(X < x, Y < y)
\]

\begin{enumerate}
\setcounter{enumi}{1}
\item Range:
\end{enumerate}

\[
0 \leq F(x, y) \leq 1
\]

\begin{enumerate}
\setcounter{enumi}{2}
\item The function is non-decreasing in both variables.
\item As $x$ and $y$ approach their maximum possible values, $F(x,y)$ approaches 1.
\item As $x$ and $y$ approach their minimum possible values, $F(x,y)$ approaches 0.
\end{enumerate}

\bigskip

\section*{End of Lecture 16}

Topics covered:

\begin{itemize}
\item Motivation for joint random variables
\item Joint PMF for discrete random variables
\item Table representation of joint distributions
\item Example with rolling two dice
\item Joint PDF for continuous random variables
\item Worked example using density $f(x,y)=4xy$
\item Joint cumulative distribution function
\item Continuous and discrete forms of joint CDF
\item Properties of the joint CDF
\end{itemize}

\end{document}